\printedpage{79}
% head: A. Grothendieck : Dérivateurs, Ch. VI, Édition des Universités de Montpellier II et Paris VII
% foot: Version 9 novembre 2001

so for $S$, $S'$ given in Cat, define
\[ \uHom(S',S)^o \longrightarrow \uHom(\HOT_W(S),\HOT_W(S')). \]
This comes down above all, for
\[ f,g : S' \rightrightarrows S \]
two arrows from $S'$ to $S$, and a homomorphism of functors
\[ u : f \longrightarrow g, \]
to describing a homomorphism of functors
\[ u^* : g^* \longrightarrow f^* \]
between $\HOT_W(S)$ and $\HOT_W(S')$. This is moreover already done, in the general setting of $(\mathcal{M},W)$, when one describes $\HOT_W(S)$ and $\HOT_W(S')$ by localization of $\mathrm{Hom}(S^o,\Cat)$ and $\mathrm{Hom}(S'^o,\Cat)$. These categories are interpreted respectively as $\uFibsc\,S$ and $\uFibsc\,S'$, hence as categories of categories over $S$ resp. $S'$, so one would like to interpret $u^*$ in terms of such categories, and while one is at it, for categories $X$ over $S$ as general as possible. We shall
\msp{page 129}
then expect to have to make assumptions, either on $X$, or on $f$ and on $g$, which ensure that the images of $\hot_{W,S}(X)$ under $f^*$ and $g^*$ are none other than $\hot_{W,S}(f^*X)$, $\hot_{W,S}(g^*X)$ respectively, where we put
\[ f^*(X) = X \times_S (S',f), \qquad g^*(X) = X \times_S (S',g). \]
So we shall probably have to place ourselves under one of the following assumptions:
\begin{enumerate}
\item[a)] $X$ is smooth over $S$.
\item[b)] $X$ is Hot-fibered over $S$.
\item[c)] $f$ and $g$ are proper (and subject to the validity of the conjecture on page 122).
\end{enumerate}
One then expects $u^*$ to be associated with a homomorphism
\[ u^* : g^*(X) \longrightarrow f^*(X) \qquad \textit{in } \Cat/S' \]
defined in terms of $u$. Is there a way of actually defining it, without assuming $X$ at least fibered over $S$? Let us take $S'=e$, so $f$, $g$ corresponding to objects $s,t\in\Ob\,S$, and $u$ to an arrow % [sic] fléche -> flèche, wrong accent
\[ u : s \longrightarrow t, \]
the inverse images of $X/S$ under $f$, $g$ are the fibers $X_s$, $X_t$, and we seek
\[ u^* : X_t \longrightarrow X_s. \]
We have no chance of finding one canonically, if we do not assume $X$ fibered over $S$! So it is not a), b) or c) that can serve us, but on%
\printedpage{80}%
% head: A. Grothendieck : Dérivateurs, Ch. VI, Édition des Universités de Montpellier II et Paris VII
% foot: Version 9 novembre 2001
\msp{page 130}%
ly the hypothesis $X\in\Ob\,\uFib\,S$. Thus one can hope to find a description of $u^*$ between $\HOT_W(S)$ and $\HOT_W(S')$ in terms of operations on the categories over $S$, $S'$, as just suggested, only by restricting oneself to ``models'' $X$ which are in $\uFib\,S$, at least. And I do not doubt indeed that this works on $\uFib\,S$ no less well than directly on $\uFibsc\,S$, in terms of $\uHom(S^o,\Cat)$.

For the essential axioms for derivators, I refer to I p. 67-68 and p. 111-112. We have seen that Der 1, \circled{Der 2} (\footnote[66]{Botched for \underline{Der 2}, I do not know whether it is true even assuming $W$ \textit{strongly} saturated, i.e. such that $\hot_W(f)$ isomorphism $\Longrightarrow f\in W$.}) are satisfied for every prederivator defined by a pair $(\mathcal{M},W\subseteq\mathrm{Fl}\,\mathcal{M})$, so it is OK here for $\HOT_W$. It follows that it is OK also for $\HOT_W^{\mathrm{lc}}$, in view of the characterization of the $\HOT_W^{\mathrm{lc}}(S)$ as subcategories of $\HOT_W(S)$. For Der 1, does this work even for arbitrary sums of categories, i.e. do we have % [sic] quelconque -> quelconques, agreement with sommes; Canulé (footnote 66) is verified authorial wording; see corrections.csv
\[ \HOT_W(\coprod_i S_i) \xrightarrow{\sim} \prod_i \HOT_W(S_i)\,? \]
We have in any case
\[ \uFib(\overbrace{\coprod\limits_i S_i}^{S}) \xrightarrow{\sim} \prod_i \uFib(S_i), \]
and the localizer $W_S^{\mathrm{f}}$ in $\uFib(S)$ is
\msp{page 131}
indeed identified with the product of the localizers $W_{S_i}^{\mathrm{f}}$ in the $\uFib\,S_i$. But in the case of a perhaps infinite family of categories $\mathcal{F}_i$, equipped with localizers $W_i$, if we put
\[ \mathcal{F} = \prod\mathcal{F}_i, \qquad W = \prod W_i \subseteq \mathrm{Fl}(\mathcal{F}), \]
is the canonical functor
\[ \mathcal{F}W^{-1} \longrightarrow \prod_i(\mathcal{F}_iW_i^{-1}) \]
always an equivalence? Even in \texttt{[``Even in'' crossed out]} the case $I$ finite one needs (it seems to me) a rather careful argument, either by the universal property by induction on card\,$I$ (one is reduced to the case of two functors), or by dint of categories of paths, noting that
\[ \uCh_\infty^W(\mathcal{F};(s_i),(t_i)) \xrightarrow{\sim} \prod_i \uCh_\infty^{W_i}(\mathcal{F}_i;s_i,t_i), \tag*{$(*)$} \]
and passing to the $\pi_0$, which gives
\[ \mathrm{Hom}_W((s_i),(t_i)) \xrightarrow{\sim} \prod_i \mathrm{Hom}_{W_i}(s_i,t_i). \]
But formula $(*)$, visibly, only works for $I$ finite. In the general case, one finds the subcategory $\prod_i^!$ of $\prod_i$ formed by the systems $(c_i)_{i\in I}$ of paths which are of bounded length.

So the question of the validity of Der 1 without a finiteness assumption on the indexing set $I$
\printedpage{81}
% head: A. Grothendieck : Dérivateurs, Ch. VI, Édition des Universités de Montpellier II et Paris VII
% foot: Version 9 novembre 2001
\msp{page 132}
remains open for the moment. (I suspect that it is true, but one would need, to prove it, to have in hand a description of the arrows in $\HOT_W(S)$ by paths of bounded length, say of length two or three. This is an interesting question, but one that I reserve for later, if need forces me to it $\ldots\,$)

More important is the question of the validity of Der 3 and of its variants (I, p. 111-112). The minimum one asks of a derivator $\mathbf{D}$ is that if $S'\longrightarrow S$ in Diag is a Hot-equivalence, it should be a $\mathbf{D}$-equivalence, i.e. (in the present case) that the functor
\[ f^* : \HOT_W(S) \longrightarrow \HOT_W(S') \]
be fully faithful on the subcategory of \textit{constant} objects in $\HOT_W(S)$. This condition becomes more tangible when one already has either the cohomology functors $H^\bullet_{\mathbf{D}}(S,-)$, $H^\bullet_{\mathbf{D}}(S',-)$, or the homology $H_\bullet^{\mathbf{D}}(S,-)$, $H_\bullet^{\mathbf{D}}(S',-)$. In terms of the latter, it means that
\[ p_!p^*(a) \xrightarrow{\sim} p'_!p'^*(a) \qquad \text{\texttt{[is an]} isomorphism,} \]
\msp{page 133}
for every $a$ in $\mathcal{A}=\mathbf{D}(e)$. (Where $p:S\longrightarrow e$, $p':S'\longrightarrow e$ are the structural morphisms to the one-point category.) We have seen the general heuristic formula
\[ p_!\underbrace{p^*(a)}_{a_S} = \hot(S)\times a, \]
where $\hot\,S$ is the class of $S\in\Ob\,\Cat$ in $\Hot=(\Cat)W_\infty^{-1}$, and where the product $\times$ is an external operation of Hot on $\mathcal{A}=\mathbf{D}(e)$
\[ \Hot\times\mathcal{A} \longrightarrow \mathcal{A}\ . \]
In the present case, $\mathcal{A}=\Hot$, and the operation
\[ \Hot\times\Hot \longrightarrow \Hot \]
should be none other than the cartesian product functor, which is deduced from the similar functor in Cat. When one has this formalism, the desired result for an $f:S'\longrightarrow S$ simply reduces to saying that if $f$ is in $W$, then so is $f\times\mathrm{id}_X$ for every $X$ in Cat (taking $X$ such that $\hot(X)\simeq a$). Now here this is a well-known fact, resulting from the axioms imposed on $W$.

So I have the impression that I have no chance of proving Der 3a$'$ (p. 111), nor even Der 3a) (the case where $\mathrm{S}=\mathrm{e}$), without having a mastery of the formalism of the $p_!$. And I have
\printedpage{82}
% head: A. Grothendieck : Dérivateurs, Ch. VI, Édition des Universités de Montpellier II et Paris VII
% foot: Version 9 novembre 2001
\msp{page 134}
good hope that this mastery will be obtained at little cost, whereas the question of the $p_*$ seems much more delicate. % [sic] cetter -> cette, misspelling

I am therefore going to turn first to the axiom Der 4$'$ of existence of
\[ f_! : \HOT_W(S') \longrightarrow \HOT_W(S), \]
as a left adjoint to $f^*$. I am going to try to construct ad-hoc a functor $f_!$, then the two adjunction arrows
\[ \begin{aligned} \mathrm{id} &\quad\longrightarrow\quad f^*f_! \\ \mathrm{id} &\quad\longleftarrow\quad f_!f^*, \end{aligned} \]
and finally, to check the two relevant compatibilities.

The ``obvious'' geometric idea is, for $X'\in\Cat/S'$, representing $\xi'\in\Ob\,\HOT_W(S')$, to define $f_!(\xi')$ by
\[ f_!(\hot_{S',W}(X')) = \hot_{S,W}(X'), \]
where on the right-hand side, $X'$ is regarded as an object of $\Cat/S$, thanks to $f:S'\longrightarrow S$ and to the functor
\[ f_! : \Cat/S' \longrightarrow \Cat/S \]
which it induces. The only problem is to see whether this functor is compatible with the
\msp{page 135}
localizers $W_{S'}^{\mathrm{d}}$, $W_S^{\mathrm{d}}$, and if not, whether it is so at least on the subcategory $\uFib\,S'$ or $\uFibsc\,S'$, which suffices to compute $\HOT_W(S')$ by localization. We must therefore see whether an $X'_1\longrightarrow X'_2$ in $\Cat/S'$, which is in $W_{S'}^{\mathrm{d}}$, i.e. which is in $W$ and remains so under every base change of the form $s'\backslash S'\longrightarrow S'$, more generally for every local co-isomorphism, is also in $W_S^{\mathrm{d}}$, i.e. whether base changes of the form $s\backslash S\longrightarrow S$ turn it into $W$-equivalences (coarse). Now local co-isomorphisms are stable under \textit{every base change} (cf. I, p. 55 in the dual case of local isomorphisms), so we win!
\[
\begin{tikzcd}
\overline{S'} = s\backslash S' \arrow[r] \arrow[d] & S' \arrow[d] \\
\overline{S} = s\backslash S \arrow[r] & S
\end{tikzcd}
\quad .
\]
I am going to check it again all the same. If
\[
\begin{tikzcd}
\overline{S'} \arrow[r] \arrow[d] & S' \arrow[d] \\
\overline{S} \arrow[r] & S
\end{tikzcd}
\]
\printedpage{83}
% head: A. Grothendieck : Dérivateurs, Ch. VI, Édition des Universités de Montpellier II et Paris VII
% foot: Version 9 novembre 2001
is a cartesian square in Cat, we said in loc. cit., essentially, that for $\overline{s'}$ in $\overline{S'}$, having images $s'$, $\overline{s}$, $s$ in $S'$, $\overline{S}$, $S$ respectively, the square of the colocalized categories
\[
\begin{tikzcd}
\overline{s'}\backslash\overline{S'} \arrow[r] \arrow[d] & s'\backslash S' \arrow[d] \\
\overline{s}\backslash\overline{S} \arrow[r] & s\backslash S
\end{tikzcd}
\]
is again
\msp{page 136}
cartesian. I have just rechecked it again this very instant, for the objects at least. To say that $\overline{S}\longrightarrow S$ is a local isomorphism, means that the $\overline{s}\backslash\overline{S}\longrightarrow s\backslash S$ are \texttt{[some]} isomorphisms. But the preceding cartesian square shows that then the $\overline{s'}\backslash\overline{S'}\longrightarrow s'\backslash S'$ are also \texttt{[some]} isomorphisms, we win! % [sic] le de revérifier -> de le revérifier, word order

\textbf{Theorem 1.} \textit{For every arrow $f:S'\longrightarrow S$ in} Cat\textit{, the induced functor}
\[ f_! : \Cat/S' \longrightarrow \Cat/S \]
\textit{is compatible with the localizers $(W_{S'}^{\mathrm{d}},W_S^{\mathrm{d}})$, hence dually with the localizers $(W_{S'}^{\mathrm{g}},W_S^{\mathrm{g}})$, and trivially with the localizers $(W_{S'},W_S)$. It therefore induces a commutative diagram of functors on the localized categories}
\[
\begin{tikzcd}
\HOT_W(S') \arrow[r, "\rho_{S'}"] \arrow[d, "f_!"] & \HOT_W^{\mathrm{lc}}(S') \arrow[d, "f_!^{\mathrm{lc}}"] & \HOT_W^o(S') \arrow[l, "\sigma_{S'}"'] \arrow[d, "f_!^o"'] \\
\HOT_W(S) \arrow[r, "\rho_S"] & \HOT_W^{\mathrm{lc}}(S) & \HOT_W^o(S). \arrow[l, "\sigma_S"']
\end{tikzcd}
\]

\textbf{Definition 1.} \textit{We have thus defined the vertical functors $f_!$, $f_!^{\mathrm{lc}}$, $f_!^o=(f^o)_!$, called homological direct image functors, or integration functors for $f$ (and relative homotopy types).}

\msp{page 137}
We have obtained at the same time, for free, the compatibility of these functors with the functors \texttt{[of]} localization $\rho_{S'}$, $\sigma_{S'}$ for $\HOT(S')$, and $\rho_S$, $\sigma_S$ for $\HOT(S)$. On the other hand, one will not expect them to commute with the fully faithful inclusion functors of the $\HOT_W^{\mathrm{lc}}$ into $\HOT_W$. Indeed, if one does not assume $f$ in \underline{Fibhot} (\footnote[67]{This case must be looked at carefully later.}), it is quite immoral to expect that $f_!$ in $\HOT_W(S')$ should transform locally constant types into likewise. The situation is therefore exactly the reverse of that for the functor $f^*$, which was compatible with the inclusion functors $i_S$, $i_{S'}$ and $j_S$, $j_{S'}$ - this was so trivial that there was no question even of saying it, let alone making a proposition of it! The situation is